\documentclass[11pt,a4paper]{article}
\usepackage[spanish,es-noquoting]{babel}
\usepackage[top=2cm, bottom=2cm, left=2.2cm, right=2.2cm]{geometry}
\usepackage{fontspec}
\usepackage{amsmath,amssymb}
\usepackage{enumitem}
\usepackage{tcolorbox}
\usepackage{microtype}

\setmainfont{Calibri}
\linespread{1.15}

\tcbset{
    colback=white,
    colframe=black,
    boxrule=0.8pt,
    arc=2mm,
    left=6pt, right=6pt, top=4pt, bottom=4pt
}

\begin{document}
\pagestyle{empty}

\begin{center}
    {\LARGE \textbf{Solucionario Detallado: MCU y Tiro Parabólico}}\\[4pt]
    {\large \textbf{Física 1.º de Bachillerato | Preparación para Cinemática Vectorial}}
\end{center}

\vspace{0.2cm}
\noindent\textbf{Resolución Ejercicio 1 (MCU y Aceleración Centrípeta):}
\begin{itemize}[leftmargin=1.5em, itemsep=3pt]
    \item \textbf{a) Período y frecuencia:}
    \[
    T = \frac{2\pi}{\omega} = \frac{2\pi}{4} = \frac{\pi}{2} \approx \mathbf{1{,}57\text{ s}} \qquad f = \frac{1}{T} = \frac{2}{\pi} \approx \mathbf{0{,}64\text{ Hz}}
    \]
    \item \textbf{b) Velocidad lineal ($v$):}
    \[
    v = \omega \cdot R = 4\text{ rad/s} \times 50\text{ m} = \mathbf{200\text{ m/s}}
    \]
    \item \textbf{c) Aceleración centrípeta o normal ($a_n$):}
    \[
    a_n = \frac{v^2}{R} = \omega^2 R = 4^2 \times 50 = \mathbf{800\text{ m/s}^2}
    \]
    \item \textbf{d) Radio de curvatura ($\rho$):}
    \[
    \rho = \frac{v^2}{a_n} = \frac{200^2}{800} = \frac{40000}{800} = \mathbf{50\text{ m}} = R
    \]
    Como la trayectoria es una circunferencia de radio $R$, su radio de curvatura es exactamente $R$ en todos los puntos.
\end{itemize}

\vspace{0.4cm}
\noindent\textbf{Resolución Ejercicio 2 (Componentes de la Aceleración):}
\begin{itemize}[leftmargin=1.5em, itemsep=3pt]
    \item \textbf{a) Aceleración normal ($a_n$):}
    \[
    a_n = \frac{v^2}{R} = \frac{20^2}{80} = \frac{400}{80} = \mathbf{5\text{ m/s}^2}
    \]
    \item \textbf{b) Aceleración total ($a$):}
    Como $\vec{a}_t \perp \vec{a}_n$:
    \[
    a = \sqrt{a_t^2 + a_n^2} = \sqrt{3^2 + 5^2} = \sqrt{9 + 25} = \sqrt{34} \approx \mathbf{5{,}83\text{ m/s}^2}
    \]
    \item \textbf{c) Radio de curvatura ($\rho$):}
    \[
    \rho = \frac{v^2}{a_n} = \frac{20^2}{5} = \frac{400}{5} = \mathbf{80\text{ m}} = R
    \]
    La aceleración tangencial solo incrementa la rapidez, mientras que la curvatura de la trayectoria está determinada únicamente por la aceleración normal que desvía la trayectoria.
\end{itemize}

\vspace{0.4cm}
\noindent\textbf{Resolución Ejercicio 3 (Tiro Parabólico y Trayectoria):}
\begin{itemize}[leftmargin=1.5em, itemsep=3pt]
    \item \textbf{a) Ecuaciones vectoriales:}
    \[
    \vec{r}(t) = (30t)\,\vec{i} + (45 - 4{,}9t^2)\,\vec{j}\text{ (m)} \qquad \vec{v}(t) = 30\,\vec{i} - 9{,}8t\,\vec{j}\text{ (m/s)}
    \]
    \item \textbf{b) Ecuación de la trayectoria $y(x)$:}
    \[
    x = 30t \implies t = \frac{x}{30} \implies y(x) = 45 - 4{,}9\left(\frac{x}{30}\right)^2 = \mathbf{45 - \frac{49}{9000}x^2 \approx 45 - 0{,}00544\,x^2}
    \]
    \item \textbf{c) Tiempo de vuelo y alcance horizontal:}
    \[
    y = 0 \implies 45 - 4{,}9t^2 = 0 \implies t = \sqrt{\frac{45}{4{,}9}} \approx \mathbf{3{,}03\text{ s}}
    \]
    \[
    x_{\text{máx}} = 30 \times 3{,}0305 = \mathbf{90{,}91\text{ m}}
    \]
    \item \textbf{d) Velocidad de impacto:}
    \[
    \vec{v}(3{,}03) = 30\,\vec{i} - 9{,}8(3{,}0305)\,\vec{j} = \mathbf{30\,\vec{i} - 29{,}70\,\vec{j}\text{ m/s}}
    \]
    \[
    v = \sqrt{30^2 + (-29{,}70)^2} = \sqrt{900 + 882} = \sqrt{1782} \approx \mathbf{42{,}21\text{ m/s}}
    \]
\end{itemize}

\newpage

\noindent\textbf{Resolución Ejercicio 4 (Cinemática Vectorial, $a_t$, $a_n$ y Curvatura):}
\begin{itemize}[leftmargin=1.5em, itemsep=3pt]
    \item \textbf{a) Vectores velocidad y aceleración:}
    \[
    \vec{v}(t) = \frac{d\vec{r}}{dt} = \mathbf{20\,\vec{i} + (15 - 10t)\,\vec{j}\text{ (m/s)}} \qquad \vec{a}(t) = \frac{d\vec{v}}{dt} = \mathbf{-10\,\vec{j}\text{ (m/s}^2\text{)}}
    \]
    \item \textbf{b) Trayectoria y altura máxima:}
    \[
    x = 20t \implies t = \frac{x}{20} \implies y(x) = 15\left(\frac{x}{20}\right) - 5\left(\frac{x}{20}\right)^2 = \mathbf{\frac{3}{4}x - \frac{1}{80}x^2}
    \]
    Altura máxima ($v_y = 0 \implies 15 - 10t = 0 \implies t = 1{,}5\text{ s}$):
    \[
    y_{\text{máx}} = y(1{,}5) = 15(1{,}5) - 5(1{,}5^2) = 22{,}5 - 11{,}25 = \mathbf{11{,}25\text{ m}}
    \]
    \item \textbf{c) En el vértice de la parábola ($t = 1{,}5\text{ s}$):}
    \begin{itemize}[leftmargin=1.5em, itemsep=2pt]
        \item $\vec{v}(1{,}5) = 20\,\vec{i} + 0\,\vec{j}\text{ m/s} \implies v = \mathbf{20\text{ m/s}}$ (horizontal).
        \item $\vec{a} = -10\,\vec{j}\text{ m/s}^2$ (vertical hacia abajo, perpendicular a $\vec{v}$).
        \item $a_t = \frac{\vec{v} \cdot \vec{a}}{v} = \frac{(20)(0) + (0)(-10)}{20} = \mathbf{0\text{ m/s}^2}$.
        \item $a_n = |\vec{a}| = \mathbf{10\text{ m/s}^2}$.
        \item Radio de curvatura: $\rho = \frac{v^2}{a_n} = \frac{20^2}{10} = \frac{400}{10} = \mathbf{40\text{ m}}$.
    \end{itemize}
    \item \textbf{d) En el instante $t = 1\text{ s}$:}
    \begin{itemize}[leftmargin=1.5em, itemsep=2pt]
        \item Vector velocidad: $\vec{v}(1) = 20\,\vec{i} + [15 - 10(1)]\,\vec{j} = \mathbf{20\,\vec{i} + 5\,\vec{j}\text{ m/s}}$.
        \item Módulo: $v(1) = \sqrt{20^2 + 5^2} = \sqrt{425} = 5\sqrt{17} \approx \mathbf{20{,}62\text{ m/s}}$.
        \item Aceleración tangencial ($a_t$):
        \[
        a_t = \frac{\vec{v} \cdot \vec{a}}{|\vec{v}|} = \frac{(20)(0) + (5)(-10)}{\sqrt{425}} = \frac{-50}{5\sqrt{17}} = -\frac{10}{\sqrt{17}} \approx \mathbf{-2{,}43\text{ m/s}^2}
        \]
        *(Signo negativo porque el cuerpo asciende frenándose)*.
        \item Aceleración normal ($a_n$):
        \[
        a_n = \sqrt{a^2 - a_t^2} = \sqrt{10^2 - \left(-\frac{10}{\sqrt{17}}\right)^2} = \sqrt{100 - \frac{100}{17}} = \sqrt{\frac{1600}{17}} = \frac{40}{\sqrt{17}} \approx \mathbf{9{,}70\text{ m/s}^2}
        \]
        \item Radio de curvatura ($\rho$):
        \[
        \rho = \frac{v^2}{a_n} = \frac{425}{\frac{40}{\sqrt{17}}} = \frac{425\sqrt{17}}{40} = \frac{85\sqrt{17}}{8} \approx \mathbf{43{,}81\text{ m}}
        \]
    \end{itemize}
\end{itemize}

\end{document}
